\documentclass[12pt]{article}

\input{../../_assets/preambulo.tex}
\DeclareMathOperator{\tr}{tr}


\begin{document}

    % 1. Foto de fondo
    % 2. Título
    % 3. Encabezado Izquierdo
    % 4. Color de fondo
    % 5. Coord x del titulo
    % 6. Coord y del titulo
    % 7. Fecha

    
    \input{../../_assets/portada}
    \portadaExamen{ffccA4.jpg}{Curvas y\\Superficies\\Examen VI}{Curvas y Superficies. Examen VI}{MidnightBlue}{-8}{28}{2026}{}

    \begin{description}
        \item[Asignatura] Curvas y Superficies.
        \item[Curso Académico] 2025/26.
        \item[Grado] Doble Grado en Ingeniería Informática y Matemáticas.
        \item[Grupo] Único.
        \item[Profesor] Sebastián Montiel Gómez.
        \item[Descripción] Convocatoria Ordinaria.
        \item[Fecha] 18 de junio de 2026.
        \item[Duración] 3 horas.
    \end{description}
    \newpage


    % ------------------------------------
    
    \noindent
    Elija y responda a continuación a los ejercicios que desee, siempre que la suma de la puntuación máxima de los mismos no supere los 10 puntos.

    \begin{ejercicio}[4 puntos]
        Enuncie el Teorema de Hilbert-Liebmann y demuéstrelo a partir del Lema (o Teorema) de Hilbert.
    \end{ejercicio}

    \begin{ejercicio}[4 puntos]
        Enuncie y demuestre el Teorema Fundamental de las Curvas Planas.
    \end{ejercicio}

    \begin{ejercicio}[2.5 puntos]
        Decide razonadamente acerca de la veracidad o falsedad de las siguientes cuestiones:
        \begin{enumerate}[label=\alph*)]
            \item Hay curvas planas regulares con curvatura constante que no son arcos de circunferencias.
            \item Hay curvas planas regulares con curvatura constante negativa.
            \item Si $S$ es una superficie conexa y orientable cuyas dos curvaturas principales son nulas en todo punto, entonces es un abierto de un plano.
        \end{enumerate}
    \end{ejercicio}

    \begin{ejercicio}[2.5 puntos]
        Demuestre que una superficie conexa y compacta con curvatura de Gau\ss\ positiva en todo punto y tal que el cociente $\nicefrac{H}{K}$ es constante tiene que ser una esfera.
    \end{ejercicio}

    \begin{ejercicio}[2.5 puntos]
        Dos superficies $S_1$ y $S_2$ se dicen tangentes en un punto $p\in \mathbb{R}^3$ si $p\in S_1\cap S_2$ y $T_pS_1 = T_pS_2$. Suponga que una superficie $S$ es tangente a un plano $P$ a lo largo de los puntos de una curva $\alpha:I\to\mathbb{R}^3$ p.p.a. Demuestre que la curvatura de Gau\ss\ $K$ de $S$ satisface $K_{\alpha(s)} = 0$ para cada $s\in I$.
    \end{ejercicio}

    \begin{ejercicio}[2.5 puntos]
        Sea $\alpha:I\to\mathbb{R}^3$ una curva p.p.a. con curvatura $k$ y torsión $\tau$ positivas. Demostrar que las rectas tangentes a $\alpha$ forman un ángulo constante con una dirección fija de $\mathbb{R}^3$ si y solo si el cociente $\nicefrac{k}{\tau}$ es constante.
    \end{ejercicio}

\end{document}
